\chapter{Thu thập và tiền xử lý dữ liệu}

\section{Nguồn dữ liệu gốc}
Dữ liệu xuất phát từ file \texttt{data/diem\_thi\_thpt\_2024.csv}, thể hiện kết quả kỳ thi tốt nghiệp THPT năm 2024 ở định dạng bảng. Các cột gốc được ánh xạ sang tên nội bộ thống nhất qua từ điển \texttt{RAW\_COLUMN\_MAP} trong \texttt{kbs/config.py} (ví dụ \texttt{ngu\_van} $\rightarrow$ \texttt{van}, \texttt{dia\_li} $\rightarrow$ \texttt{dia\_ly}) \cite{docs_dataset}.

\section{Sơ đồ pipeline tiền xử lý (mức logic)}
Luồng xử lý trong \texttt{scripts/create\_data.py} có thể tóm tắt như sau (độc lập ngôn ngữ lập trình):
\begin{verbatim}
  RAW (diem_thi_thpt_2024.csv)
       | doi ten cot theo RAW_COLUMN_MAP
       | loc du 6 mon theo khoi (KHTN / KHXH)
       | gan nhan nganh_hoc (heuristic theo khoi)
       | can bang lop (undersampling)
       v
  data_khtn.csv / data_khxh.csv
\end{verbatim}
Mỗi bước đều có thể ghi log số dòng trước/sau để phục vụ kiểm toán dữ liệu trong báo cáo đồ án.

\section{Lọc khối và điều kiện đủ môn}
\texttt{scripts/create\_data.py} thực hiện các bước: đọc và đổi tên cột; lọc thí sinh thuộc khối KHTN nếu có đủ ba môn tự chọn Lý--Hóa--Sinh cùng ba môn bắt buộc; lọc KHXH nếu có đủ Lịch sử--Địa lý--GDCD. Các hàng thiếu giá trị ở môn bắt buộc bị loại. Điểm nằm ngoài $[0,10]$ được kẹp về biên theo chính sách mô tả trong tài liệu dữ liệu.

\section{Heuristic gán nhãn ngành}
Nhãn \texttt{nganh\_hoc} là số nguyên trùng khóa trong \texttt{NGANH\_HOC\_MAP}. Với KHTN, hàm \texttt{assign\_major\_khtn} xây dựng từ điển điểm ưu tiên khi các điều kiện ngưỡng đồng thời thỏa (ví dụ IT khi Toán $\geq 7{,}5$ và Lý $\geq 7$); ngành được chọn là khóa có điểm ưu tiên lớn nhất, nếu không có khớp tuyệt đối thì dùng hàm tổ hợp tuyến tính fallback trên toàn bộ ngành. Khối KHXH dùng \texttt{assign\_major\_khxh} với logic tương tự cho bốn ngành kinh tế, sư phạm, luật, du lịch.

Việc gán nhãn tự động cho phép huấn luyện có giám sát quy mô lớn nhưng \textbf{không} tương đương khảo sát định hướng thực tế: báo cáo luôn ghi nhận độ lệch tiềm ẩn giữa nhãn proxy và hành vi chọn ngành của học sinh.

\section{Cân bằng lớp}
Sau khi gán nhãn, phân bố lớp thường lệch. Hàm \texttt{balance\_data} thực hiện \textbf{undersampling} ngẫu nhiên về cỡ lớp nhỏ nhất trong tập ngành hợp lệ của khối, đảm bảo mỗi lớp có cùng số mẫu trước khi ghi CSV. Điều này giúp mô hình không bị thiên lệch tần suất quá mức về một ngành phổ biến trong heuristic.

\section{Thống kê mô tả trên tập đã xử lý}
Bảng \ref{tab:khoin} và \ref{tab:khoixh} tổng hợp số mẫu, phân bố lớp và thống kê mô tả sáu đặc trưng \textbf{đo trực tiếp} từ các file \texttt{data/data\_khtn.csv} và \texttt{data/data\_khxh.csv} tại thời điểm báo cáo (máy chủ mã dự án).

\begin{table}[htbp]
\centering
\caption{Dữ liệu KHTN sau cân bằng ($N=167\,490$)}
\label{tab:khoin}
\begin{tabular}{lrrrrr}
\toprule
Ngành (mã) & Số mẫu & Tỉ lệ \% \\
\midrule
0 IT & 33\,498 & 20{,}0 \\
1 Kinh tế & 33\,498 & 20{,}0 \\
2 Y khoa & 33\,498 & 20{,}0 \\
3 Kỹ thuật & 33\,498 & 20{,}0 \\
4 Nông--Lâm--Ngư & 33\,498 & 20{,}0 \\
\midrule
\textbf{Tổng} & \textbf{167\,490} & \textbf{100} \\
\bottomrule
\end{tabular}
\end{table}

\begin{table}[htbp]
\centering
\caption{Thống kê mô tả sáu đặc trưng KHTN (cùng tập)}
\begin{tabular}{lrrrrrr}
\toprule
 & toan & van & anh & ly & hoa & sinh \\
\midrule
Trung bình & 7{,}537 & 7{,}286 & 5{,}967 & 6{,}572 & 6{,}861 & 6{,}523 \\
Đ.lệch chuẩn & 0{,}932 & 1{,}119 & 1{,}744 & 1{,}617 & 1{,}601 & 1{,}258 \\
Min / Max & 1{,}6 / 9{,}8 & 0{,}75 / 10 & 0{,}6 / 10 & 0{,}75 / 10 & 1 / 10 & 0 / 10 \\
\bottomrule
\end{tabular}
\end{table}

\begin{table}[htbp]
\centering
\caption{Dữ liệu KHXH sau cân bằng ($N=63\,848$)}
\label{tab:khoixh}
\begin{tabular}{lrr}
\toprule
Ngành (mã) & Số mẫu & Tỉ lệ \% \\
\midrule
1 Kinh tế & 15\,962 & 25{,}0 \\
5 Sư phạm & 15\,962 & 25{,}0 \\
6 Luật & 15\,962 & 25{,}0 \\
7 Du lịch & 15\,962 & 25{,}0 \\
\midrule
\textbf{Tổng} & \textbf{63\,848} & \textbf{100} \\
\bottomrule
\end{tabular}
\end{table}

\begin{table}[!htbp]
\centering
\caption{Thống kê mô tả sáu đặc trưng KHXH (cùng tập $N=63\,848$)}
\begin{tabular}{lrrrrrr}
\toprule
 & toan & van & anh & lich\_su & dia\_ly & gdcd \\
\midrule
Trung bình & 6{,}731 & 7{,}696 & 6{,}734 & 6{,}284 & 7{,}112 & 7{,}948 \\
Đ.lệch chuẩn & 1{,}356 & 1{,}305 & 1{,}787 & 1{,}233 & 1{,}148 & 1{,}178 \\
Min / Max & 1 / 9{,}6 & 0 / 10 & 0{,}4 / 10 & 0{,}5 / 10 & 0 / 10 & 0 / 10 \\
\bottomrule
\end{tabular}
\end{table}
\FloatBarrier

Nhận xét định tính: khối KHTN có điểm Toán trung bình cao hơn KHXH trong tập đã cân bằng; môn Ngoại ngữ ở KHTN có độ phân tán lớn hơn (độ lệch chuẩn $\approx 1{,}74$). Khối KHXH có điểm GDCD trung bình cao, phản ánh tổ hợp môn xã hội.

\section{Chia tập huấn luyện và kiểm tra}
Toàn bộ pipeline huấn luyện (\texttt{scripts/train\_model.py}) dùng \texttt{train\_test\_split} với \texttt{test\_size=0{,}2}, \texttt{stratify=y}, \texttt{random\_state=42}. Mọi số liệu độ chính xác trong chương 7 tuân theo cùng một lần chia này để đảm bảo có thể tái lập.


\section{Nhãn thay thế (\textit{surrogate labels}) và hệ quả đánh giá}
Heuristic gán nhãn tạo ra nhãn ``ứng viên'' $y$ không phải lựa chọn thực tế của thí sinh mà là xấp xỉ từ điểm số. Trong học máy, loại nhãn này thường gọi là nhãn giám sát yếu (\textit{weak supervision}). Hệ quả: mô hình có thể đạt độ chính xác cao trên tập kiểm tra nếu heuristic tương quan mạnh với đặc trưng; song metric \textbf{không} tương đương độ đúng trên hành vi chọn ngành thực tế. Đồ án luôn tách bạch hai khái niệm này khi diễn giải chương 7.

\section{Đạo đức và quyền riêng tư dữ liệu}
Dữ liệu điểm thi gắn với cá nhân nếu có định danh. Nguyên tắc tối thiểu hóa dữ liệu: chỉ giữ sáu môn và nhãn nội bộ; không lưu họ tên trong CSV đã xử lý của đồ án. Khi triển khai công cộng cần thông báo rõ mục đích xử lý và cơ chế xóa dữ liệu.

\section{Ma trận chất lượng dữ liệu (5V)}
\begin{description}[style=nextline,leftmargin=0cm]
  \item[Volume] Sau cân bằng, KHTN $\approx 1{,}67\times 10^5$ mẫu; đủ lớn cho RF nhưng vẫn một năm.
  \item[Velocity] Cập nhật theo kỳ thi hàng năm; pipeline batch phù hợp.
  \item[Variety] Chỉ số vô hướng; chưa có dữ liệu đa phương thức.
  \item[Veracity] Nguồn chính thức; nhãn heuristic cần kiểm chứng ngoài thực địa.
  \item[Value] Giá trị nằm ở khả năng minh họa hệ lai có giải thích.
\end{description}

\section{Kiểm định nhất quán nội bộ}
Trước khi huấn luyện, nên kiểm tra: (i) mỗi hàng có đúng sáu cột số; (ii) \texttt{nganh\_hoc} thuộc tập cho phép theo khối; (iii) không có giá trị null sau lọc. Có thể viết script assert trong CI để phát hiện lệch lược đồ CSV.

\section{Tổng kết chương}
Dữ liệu dự án có quy mô lớn, cân bằng lớp rõ ràng và tiền xử lý minh bạch trong mã nguồn. Hạn chế trọng tâm là tính proxy của nhãn; các chương sau khai thác tri thức và học máy trên cơ sở đó.
